\section{Branch messages and rooted stackability}\label{sec:messages}
For an oriented edge $v\to p$, let $B_{v\mid p}$ be the component
containing $v$ after deleting $vp$. Its root is $v$ and its external boundary
vertex is $p$. The children are the branches $B_{u\mid v}$ with
$u\sim v$, $u\ne p$. Call a branch \emph{occupied} when its restriction of
$C$ is nonzero.

Define the integer transfer function by
\begin{equation}\label{eq:transfer}
 F(x)=\begin{cases}
 2x-3,&x\le1,\\
 x/2,&x\ge2\text{ even},\\
 (x-3)/2,&x\ge3\text{ odd}.
 \end{cases}
\end{equation}
In particular, $F(2)=1$ and $F(3)=0$.
Messages take values in the disjoint union $\{\EMPTY\}\sqcup\ZZ$.
For an empty branch put $m_{v\to p}=\EMPTY$. For an occupied branch put
\begin{equation}\label{eq:message}
 m_{v\to p}=d_{v\to p}:=
 F\left(C(v)+\sum_{\substack{u\sim v,\ u\ne p\\B_{u\mid v}\text{ occupied}}}
                          d_{u\to v}\right).
\end{equation}
The recursion terminates because each child branch has fewer vertices.
Only integer messages are summed. When writing a sum over all neighbours
we use the numerical contribution $\bar d_{u\to v}$, equal to
$d_{u\to v}$ on occupied branches and $0$ on empty branches. This
convention for sums does not identify the states $\EMPTY$ and $0$.
The rooted score is
\begin{equation}\label{eq:score}
 S_r(C)=C(r)+\sum_{u\sim r}\bar d_{u\to r}.
\end{equation}

\begin{theorem}[Exact boundary invariant]\label{thm:boundary}
Let $B=B_{v\mid p}$ be occupied with message $d=d_{v\to p}$.
Work on the subtree induced by $B\cup\{p\}$, with initial configuration
$C\restriction B$ on $B$ and $q\ge0$ pebbles at $p$.
A positive stack at $p$ with $B$ empty is reachable if and only if
$q+d>0$. In that case the greatest possible final pile at $p$ is $q+d$.
\end{theorem}

We prove necessity in a form that also applies when other branches
interleave their moves at $p$.
\begin{lemma}[Boundary flux]\label{lem:flux}
Suppose a legal sequence clears $B_{v\mid p}$. Let $m_{a b}$ count its
moves from $a$ to $b$, and put
$\phi_{v\mid p}=m_{v p}-2m_{p v}$.
If the branch is occupied initially, then
\[
 \phi_{v\mid p}\le d_{v\to p},\qquad
 \phi_{v\mid p}\equiv d_{v\to p}\pmod3.
\]
If it is empty initially, then $\phi_{v\mid p}=-3k$ for some $k\ge0$.
Moves elsewhere in the tree are permitted.
\end{lemma}
\begin{proof}
The move counts satisfy, at every cleared vertex $w$,
\begin{equation}\label{eq:balance}
 C(w)+\sum_{z\sim w}m_{z w}-2\sum_{z\sim w}m_{w z}=0.
\end{equation}
For each initially occupied branch which is cleared, at least one move
crosses its boundary outwards. Indeed, consider the move after which that
branch becomes empty for the last time. An internal move leaves a pebble
inside, and an inward move also does so. Thus this move is an outward
boundary move. This observation applies to every occupied child branch,
independently of the moves outside it.

First consider an initially empty branch. Sum~\eqref{eq:balance} over
its vertices. If $I$ is the number of internal moves and
$a=m_{v p}$, $b=m_{p v}$, the result is $b-2a-I=0$.
Hence its boundary flux is
$a-2b=-3a-2I\le0$. It is divisible by three: give the boundary vertex
sign $+1$ and alternate signs along tree edges. The signed sum of the
balance equations over the branch is congruent to minus its boundary
flux modulo three, because each internal move contributes a multiple of
three. The initial and final signed sums on the empty branch vanish.
This proves the assertion for empty branches, even when they are used
temporarily.

For an occupied branch, induct on its size. By induction, each occupied
child has flux $d_i-3k_i$ with $k_i\ge0$, and each empty child has flux
$-3k_i$. Put $x=C(v)+\sum d_i$ and $K=\sum k_i$, where the first sum
runs over occupied children. At $v$, equation~\eqref{eq:balance} reads
\[
 y+b-2a=0,\qquad y=x-3K,\quad a=m_{v p}\ge1,\quad b=m_{p v}\ge0.
\]
For any integer $y$, these conditions imply
\begin{equation}\label{eq:one-vertex-flux}
 a-2b=2y-3a\le F(y),\qquad a-2b\equiv F(y)\pmod3.
\end{equation}
To see this, the least possible $a$ is
$\max(1,\lceil y/2\rceil)$; substituting this value gives precisely
the three cases of~\eqref{eq:transfer}. Any increase in $a$ lowers
the flux by three.

The transfer function satisfies
\begin{equation}\label{eq:three-loss}
 F(z-3)\le F(z),\qquad F(z)\equiv2z\pmod3.
\end{equation}
For the inequality, the difference $F(z)-F(z-3)$ is $6$ when $z\le1$,
$6,3,3$ at $z=2,3,4$, respectively, and is $0$ for odd $z\ge5$ and
$3$ for even $z\ge6$. The congruence follows from each case in
\eqref{eq:transfer}. Thus $F(x-3K)\le F(x)$ and
$F(x-3K)\equiv F(x)\pmod3$. Apply~\eqref{eq:one-vertex-flux}
to finish the induction.
\end{proof}

\begin{proof}[Proof of Theorem~\ref{thm:boundary}]
The final boundary pile in any allowed clearing sequence is $q+\phi$,
so Lemma~\ref{lem:flux} proves necessity and the upper bound.
For attainment we induct on the branch size.

We use the following scheduling observation. Suppose disjoint occupied
child branches have gains $d_i$, and a vertex initially has $a\ge0$
pebbles. If $a+\sum_i d_i>0$, process branches with positive gain first,
then those with zero gain, then those with negative gain. A positive-gain
task has positive final pile even when begun at zero. Before a zero-gain
task the pile is positive, since the total remaining gains are nonpositive
and the final total is positive. During the negative-gain tasks the pile
after each task is at least the final positive total. The inductive
boundary invariant therefore realizes every task and leaves exactly
$a+\sum_i d_i$. Tasks touch only their branch and the common boundary;
no other task's initial interior configuration is changed.

Let $x=C(v)+\sum_i d_i$ for the occupied children of $B$.
If $x\ge2$, the observation clears the children and leaves $x$ at $v$.
For $x=2k$, make $k$ moves from $v$ to $p$; the final boundary pile is
$q+k=q+F(x)$.
For $x=2k+1\ge3$, make $k$ moves from $v$ to $p$, leaving one pebble
at $v$. The assumption $q+F(x)=q+k-1>0$ implies $q+k\ge2$.
One move $p\to v$ followed by one move $v\to p$ is therefore legal;
it clears $v$ and leaves $q+k-1$ at $p$.

If $x\le1$, put $j=2-x\ge1$. The assumption
$q+F(x)=q-2j+1>0$ implies $q\ge2j$. Make $j$ moves $p\to v$.
The children's effective total at $v$ is now $x+j=2$.
The observation clears them leaving two pebbles at $v$, and one final
move $v\to p$ gives pile $q-2j+1=q+F(x)$.
These constructions also cover the case with no children and complete
the induction.
\end{proof}

\begin{theorem}[Exact rooted stackability]\label{thm:root-score}
For every configuration $C$ on a finite tree and every vertex $r$,
\[
 \StackableAt(T,C,r)\quad\Longleftrightarrow\quad S_r(C)>0.
\]
Consequently $C$ is nonstackable if and only if all its scores are
nonpositive.
\end{theorem}
\begin{proof}
If the score is positive, apply the scheduling observation to the occupied
branches incident with $r$, starting with $C(r)$ pebbles. It gives a stack
of size exactly $S_r(C)$ at $r$.
Conversely, in any sequence ending in a stack at $r$, every incident
branch is cleared. Lemma~\ref{lem:flux} bounds the flux from each occupied
branch by its message and from each empty branch by zero. The final pile
is therefore at most $S_r(C)$, and it is positive.
The final equivalence follows by quantifying over all targets.
\end{proof}

\begin{remark}
An occupied one-vertex branch containing three pebbles has message
$F(3)=0$. It can be cleared to its boundary only when the boundary has a
positive initial pile. An empty branch may be left alone even when that
pile is zero. This is why the recursive state must retain $\EMPTY$
separately from the integer zero.
\end{remark}
